\chapter{Mapeamentos Fornecedores}
\label{apendice:mapeamentos-fornecedores}

    \section{Tradutores de Entrada (Inbound)}
    \subsection{\textbf{OpenAI} Inbound Translator}

A implementação \texttt{OpenAiInboundTranslator} adapta os contratos compatíveis com a API da \textbf{OpenAI} para o domínio comum.
    
\begin{table}[H]
    \centering
    \footnotesize
    \renewcommand{\arraystretch}{1.25}
    \begin{tabularx}{\textwidth}{|>{\ttfamily\raggedright\arraybackslash}p{4cm}|Y|}
        \hline
        \textbf{Contrato Inbound (OpenAI)} & \textbf{Mapeamento para CommonRequest} \\
        \hline
        temperature, maxTokens, topP &
        Integrados na \texttt{CommonGenerationConfig}. O campo \texttt{model} é mapeado diretamente para o nível superior do \texttt{CommonRequest}. \\
        \hline
        stop &
        É convertido para uma \texttt{List<String>}, distinguindo automaticamente entre valores únicos (\texttt{Single}) e múltiplos (\texttt{Multiple}). \\
        \hline
        messages[].role &
        Convertido para os papéis do domínio comum: \texttt{SYSTEM}, \texttt{ASSISTANT}, \texttt{TOOL} e \texttt{USER}. Valores não reconhecidos são tratados como \texttt{USER}. \\
        \hline
        messages[].content \& toolCalls &
        O conteúdo textual não vazio origina instâncias de \texttt{TextPart}. As chamadas de ferramentas são convertidas em \texttt{ToolCallPart(toolCallId, functionName)}, com os argumentos encapsulados num mapa sob a chave \texttt{"raw"}; se o identificador vier omisso, é gerado no formato \texttt{"openai-tool-call-\$index"}. Quando o papel é \texttt{TOOL} e existem \texttt{toolCallId} e conteúdo, é gerado um \texttt{ToolResultPart}, convertendo o conteúdo textual para valor JSON. \\
        \hline
        tools[].function &
        Convertido para \texttt{CommonTool(name, description, parametersSchema)}, extraindo as propriedades definidas no schema da ferramenta. Se a descrição estiver ausente, é usado o nome da função; se os parâmetros estiverem ausentes, o schema fica vazio. \\
        \hline
        toolChoice &
        Converte os valores \texttt{auto}, \texttt{none} e \texttt{required} para os respetivos tipos do domínio comum (\texttt{ToolChoice.Auto}, \texttt{ToolChoice.None} e \texttt{ToolChoice.Required}), podendo também representar a seleção explícita de uma função específica. \\
        \hline
        responseFormat &
        Verifica se o tipo corresponde a \texttt{"json\_object"} ou \texttt{"json\_schema"}, ativando a opção \texttt{jsonResponse}. \\
        \hline
        \textit{stream}, frequencyPenalty, presencePenalty, n, seed, user, logitBias, logProbs, topLogProbs, responseFormat &
        Armazenados na \texttt{TypedMap} (\texttt{providerOptions}) do \texttt{CommonRequest}, preservando opções específicas do fornecedor que não possuem representação no domínio comum. \\
        \hline
    \end{tabularx}
    \caption{Mapeamento Inbound: \textbf{OpenAI} Request para o Domínio Comum}
    \label{tab:openai-inbound-mapping}
\end{table}


No sentido inverso, a tradução da resposta converte os objetos do domínio comum para a estrutura nativa da API da \textbf{OpenAI}.

\begin{table}[H]
    \centering
    \footnotesize
    \renewcommand{\arraystretch}{1.25}
    \begin{tabularx}{\textwidth}{|>{\ttfamily\raggedright\arraybackslash}p{4cm}|Y|}
        \hline
        \textbf{Domínio Comum} & \textbf{Mapeamento para Contrato Resposta (OpenAI)} \\
        \hline
        id \& created &
        Utiliza o \texttt{id} nativo ou gera um identificador no formato \texttt{chatcmpl-UUID}. O instante temporal é recuperado das \texttt{providerOptions} ou gerado pelo \texttt{Clock.System}; o \texttt{systemFingerprint}, quando existe, também é recuperado das \texttt{providerOptions}. \\
        \hline
        choices[].message.role &
        Convertido de \texttt{CommonRole} para os papéis nativos da \textbf{OpenAI}: \texttt{system}, \texttt{user}, \texttt{assistant} ou \texttt{tool}. \\
        \hline
        choices[].message.content &
        O conteúdo é obtido através de \texttt{firstRenderableContentOrNull()}, juntando os \texttt{TextPart} com quebras de linha. Se não existir texto, recorre ao primeiro \texttt{RefusalPart} ou à representação textual do primeiro \texttt{JsonPart}. \\
        \hline
        choices[].message
        .toolCalls &
        Um \texttt{ToolCallPart} é convertido para um \texttt{OpenAiToolCallOutput}, formatando os \texttt{argumentsJson} numa única representação textual. \\
        \hline
        choices[].finishReason &
        Converte \texttt{STOP} para \texttt{stop}, \texttt{LENGTH} para \texttt{length}, \texttt{TOOL\_CALL} para \texttt{tool\_calls} e \texttt{CONTENT\_FILTER} para \texttt{content\_filter}. Valores \texttt{OTHER} ou nulos não são emitidos. \\
        \hline
        usage &
        Mapeia \texttt{inputTokens}, \texttt{outputTokens} e \texttt{totalTokens} para \texttt{promptTokens}, \texttt{completionTokens} e \texttt{totalTokens}, respetivamente. \\
        \hline
    \end{tabularx}
    \caption{Mapeamento Inbound: Domínio para \textbf{OpenAI} Response Unária}
\end{table}


Para respostas em \textit{streaming}, os eventos internos do domínio são convertidos para os \textit{chunks} \textbf{SSE} utilizados pela \textbf{OpenAI}.

\begin{table}[H]
    \centering
    \footnotesize
    \renewcommand{\arraystretch}{1.25}
    \begin{tabularx}{\textwidth}{|>{\raggedright\arraybackslash}p{4cm}|Y|}
        \hline
        \textbf{Common Response Event} & \textbf{Mapeamento para Stream Chunk (OpenAI)} \\
        \hline
        \texttt{ResponseStarted} & Emite Chunk inicial. \texttt{choices} ficam vazias. \\
        \hline
        \texttt{ChoiceStarted} & Emite Chunk com índice da \textit{choice}. O \texttt{delta} contém apenas a definição do \texttt{role}; se o papel vier ausente, assume \texttt{assistant}. \\
        \hline
        \texttt{TextDeltaEvent} & Emite Chunk com \texttt{delta.content} preenchido pelo texto. \\
        \hline
        \texttt{ToolCallStartedEvent} & Emite Chunk inicializando o delta de \texttt{toolCalls}: define o \texttt{id}, a \texttt{type="function"} e o nome da função (\texttt{arguments=""}). \\
        \hline
        \texttt{ToolCallArguments
        DeltaEvent} & Emite Chunk contendo fragmentos incrementais em \texttt{function.arguments}. O nome da função é omitido neste pacote. \\
        \hline
        \texttt{ChoiceFinished} & Emite Chunk apenas com a tradução do \texttt{finishReason}; o \texttt{delta} não é preenchido. \\
        \hline
        \texttt{UsageReported} & Emite Chunk especial sem \textit{choices}, transportando o objeto global \texttt{usage}. \\
        \hline
        \texttt{ResponseCompleted} & Emite Chunk sem \textit{choices}, sinalizando a conclusão bem-sucedida do \textit{stream}. \\
        \hline
        \texttt{ResponseErrored} & Emite Chunk injetando a mensagem de erro como texto no \texttt{delta.content}. \\
        \hline
    \end{tabularx}
\caption{Mapeamento Inbound: Domínio para \textbf{OpenAI} Streaming (\texttt{fromDomainEvent})}
\end{table}

\subsection{\textbf{Anthropic} Inbound Translator}

O \texttt{AnthropicInboundTranslator} adapta os contratos compatíveis com a \textit{Messages API} da \textbf{Anthropic} para o domínio comum.

\begin{table}[H]
    \centering
    \footnotesize
    \renewcommand{\arraystretch}{1.25}
    \begin{tabularx}{\textwidth}{|>{\ttfamily\raggedright\arraybackslash}p{4.5cm}|Y|}
        \hline
        \textbf{Contrato Inbound (Anthropic)} & \textbf{Mapeamento para CommonRequest} \\
        \hline
        temperature, maxTokens, topP, stopSequences & Agrupados na \texttt{CommonGenerationConfig}. O campo \texttt{model}, por sua vez, fica mapeado logo na raiz do \texttt{CommonRequest}. \\
        \hline
        system & O conteúdo de sistema é convertido para \texttt{SystemPrompt}. Quando chega como \texttt{RawText}, o texto é usado diretamente; quando chega como lista de blocos, apenas os blocos \texttt{Text} são unidos com \texttt{joinToString()}. Se o resultado estiver em branco, é ignorado. \\
        \hline
        messages[].role & Convertido para os enums de domínio: \texttt{SYSTEM}, \texttt{ASSISTANT}, \texttt{TOOL} ou \texttt{USER}. Valores não reconhecidos são tratados como \texttt{USER}. \\
        \hline
        messages[].content & O \texttt{Text} passa a \texttt{TextPart}. O \texttt{ToolUse} dá origem a um \texttt{ToolCallPart(id, name, input)}, gerando um ID automático (\texttt{"anthropic-tool-use-\$index"}) se for necessário. O \texttt{ToolResult} é encapsulado num \texttt{ToolResultPart}. \textbf{Nota}: Blocos de \texttt{Thinking} (raciocínio interno) que venham no input são ignorados e retornam \texttt{null}. \\
        \hline
        tools & A lista é mapeada para instâncias de \texttt{CommonTool}. O \texttt{inputSchema} original é convertido para \texttt{JsonObject} e guardado no campo \texttt{parametersSchema}; se não for um objeto JSON, o schema fica vazio. \\
        \hline
        toolChoice & Traduzido diretamente: \texttt{auto} para \texttt{Auto}, \texttt{none} para \texttt{None}, \texttt{any} para \texttt{Required}, e os casos de \texttt{tool} passam a \texttt{Specific(name)}. \\
        \hline
        \normalfont\textit{(Constantes da Integração)} & O \texttt{provider} fica definido estaticamente como \texttt{Provider.ANTHROPIC} e o \texttt{jsonResponse} assume sempre o valor \texttt{false}. \\
        \hline
        \textit{stream}, topK, stopToken, thinking, metadata & Por serem opções específicas e mais avançadas, são guardadas no \texttt{TypedMap} (\texttt{providerOptions}) dentro do \texttt{CommonRequest}. \\
        \hline
    \end{tabularx}
    \caption{Mapeamento Inbound: Request da \textbf{Anthropic} para o Domínio}
\end{table}

No sentido inverso, a resposta do domínio comum é convertida para o formato \texttt{Anthropic\allowbreak{}Message\allowbreak{}Response}.


\begin{table}[H]
    \centering
    \footnotesize
    \renewcommand{\arraystretch}{1.25}
    \begin{tabularx}{\textwidth}{|>{\ttfamily\raggedright\arraybackslash}p{4cm}|Y|}
        \hline
        \textbf{Domínio Comum} & \textbf{Mapeamento para Contrato Resposta (Anthropic)} \\
        \hline
        id \& model & Utiliza diretamente os valores de \texttt{id} e \texttt{model} da \texttt{CommonResponse}. A resposta nativa é construída a partir da primeira \textit{choice}; se não existir, o conteúdo fica vazio e o papel assume \texttt{assistant}. \\
        \hline
        choices[].message.role & Convertido de \texttt{CommonRole} para o papel nativo. \texttt{USER} passa a \texttt{user}; \texttt{SYSTEM}, \texttt{ASSISTANT} e \texttt{TOOL} passam a \texttt{assistant}. \\
        \hline
        choices[].message.content & O conteúdo é convertido preservando os elementos nativos: \texttt{TextPart} gera \texttt{Text}, \texttt{ToolCallPart} origina \texttt{ToolUse} preservando a hierarquia do JSON interno, \texttt{JsonPart} é serializado como texto JSON, e os \texttt{RefusalPart} são omitidos do mapeamento. \\
        \hline
        choices[].finishReason & Convertido para \texttt{end\_turn}, \texttt{max\_tokens}, \texttt{tool\_use} e \texttt{stop\_sequence}. Valores \texttt{OTHER} ou nulos não são emitidos. \\
        \hline
        usage & Mapeado para o correspondente \texttt{AnthropicUsage} (\texttt{inputTokens} e \texttt{outputTokens}). \\
        \hline
    \end{tabularx}
\caption{Mapeamento Inbound: Domínio para \textbf{Anthropic} Message Response Unária}
\end{table}

Em modo \textit{streaming}, os eventos do domínio são traduzidos para os eventos \textbf{SSE} definidos pela \textbf{Anthropic}.
\begin{table}[H]
    \centering
    \footnotesize
    \renewcommand{\arraystretch}{1.25}
    \begin{tabularx}{\textwidth}{|>{\raggedright\arraybackslash}p{4.5cm}|Y|}
        \hline
        \textbf{Common Response Event} & \textbf{Mapeamento para Evento SSE (Anthropic)} \\
        \hline
        \texttt{ResponseStarted} & Retorna um evento \texttt{MessageStart} com o \texttt{id} e o \texttt{model} do evento de domínio, assumindo o papel de \texttt{role="assistant"}. \\
        \hline
        \texttt{ChoiceStarted} & Emite um \texttt{ContentBlockStart} com o índice da \textit{choice} correspondente, indicando o início de um bloco de \texttt{Text}. \\
        \hline
        \texttt{TextDeltaEvent} & O fragmento de texto (\texttt{text}) é repassado de forma direta para o evento \textbf{SSE} \texttt{ContentBlockDelta(TextDelta)}. \\
        \hline
        \texttt{ToolCallStartedEvent} & Inicia o bloco com um \texttt{ContentBlockStart(ToolUse)}, mantendo o \texttt{id} e o \texttt{name} originais, e inicializando a chamada com um \texttt{JsonObject} vazio. \\
        \hline
        \texttt{ToolCallArgumentsDeltaEvent}& Encaminha os fragmentos parciais dos argumentos recebidos para um evento \textbf{SSE} do tipo \texttt{ContentBlockDelta(InputJsonDelta)}. \\
        \hline
        \texttt{ChoiceFinished} & Gera um evento \texttt{MessageDelta} que inclui o motivo de paragem (\texttt{stopReason}) associado a esta \textit{choice}. \\
        \hline
        \texttt{UsageReported} & Gera um evento \texttt{MessageDelta} contendo a contagem atualizada de \textit{tokens} processados. \\
        \hline
        \texttt{ResponseCompleted} & Emite o evento \texttt{MessageStop}, sinalizando o fecho definitivo da \textit{stream}. \\
        \hline
        \texttt{ResponseErrored} & Interrompe a \textit{stream} com um evento de \texttt{Error}. Faz-se a distinção entre falhas temporárias (\texttt{overloaded\_error}) e erros gerais da API (\texttt{api\_error}). \\
        \hline
    \end{tabularx}
    \caption{Mapeamento Inbound: Eventos \textbf{SSE} \textbf{Anthropic} (\texttt{fromDomainEvent})}
\end{table}

    \subsection{\textbf{Gemini} Inbound Translator}

O \texttt{GeminiInboundTranslator} adapta os contratos compatíveis com a API \textbf{Gemini} para o domínio comum.

\begin{table}[H]
    \centering
    \footnotesize
    \renewcommand{\arraystretch}{1.25}
    \begin{tabularx}{\textwidth}{|>{\ttfamily\raggedright\arraybackslash}p{4.5cm}|Y|}
        \hline
        \textbf{Contrato Inbound (Gemini)} & \textbf{Mapeamento para CommonRequest} \\
        \hline
        generationConfig & As propriedades base (\texttt{temperature}, \texttt{topP}, \texttt{stopSequences}) são extraídas diretamente para a \texttt{CommonGenerationConfig}. \\
        \hline
        model & O identificador do modelo é lido do contrato recebido. Caso não seja fornecido, assume-se a constante \texttt{"NO MODEL"} como \textit{fallback}. \\
        \hline
        systemInstruction & Convertido para um \texttt{SystemPrompt}, agrupando as \texttt{parts} de texto não vazias num único bloco contínuo, separadas por quebras de linha. Se o resultado ficar vazio, o \texttt{SystemPrompt} não é criado. \\
        \hline
        contents[].role & Adapta os papéis nativos para o domínio: \texttt{system} passa a \texttt{SYSTEM}, \texttt{model} e \texttt{assistant} passam a \texttt{ASSISTANT}, \texttt{tool} e \texttt{function} passam a \texttt{TOOL}, enquanto \texttt{user} se mantém como \texttt{USER}. Valores não reconhecidos são tratados como \texttt{USER}. \\
        \hline
        contents[].parts & O conteúdo é transformado de acordo com o seu tipo: \texttt{text} gera um \texttt{TextPart}, \texttt{functionCall} origina um \texttt{ToolCallPart} com identificador local \texttt{"gemini-tool-call-\$index"}, e \texttt{functionResponse} converte-se em \texttt{ToolResultPart}. \\
        \hline
        tools[].functionDeclarations & A hierarquia aninhada (\texttt{tools -> functionDeclarations}) é simplificada (\textit{flattened}) para criar uma lista única de ferramentas (\texttt{List<CommonTool>}). \\
        \hline
        toolConfig\newline.functionCallingConfig & Se existirem \texttt{allowedFunctionNames}, o tradutor cria uma escolha específica. Caso contrário, adapta o modo de execução: \texttt{auto} passa a \texttt{Auto}, \texttt{none} a \texttt{None}, e \texttt{any} ou \texttt{required} a \texttt{Required}. \\
        \hline
        generationConfig\newline.responseMimeType & A \textit{flag} \texttt{jsonResponse} assume o valor \texttt{true} caso o tipo MIME seja \texttt{application/json} ou se existir um \texttt{responseJsonSchema} explicitamente definido. \\
        \hline
        topK, thinkingConfig, responseMimeType, responseJsonSchema & Por se tratarem de configurações específicas do \textbf{Gemini}, são guardadas no \texttt{TypedMap} (\texttt{providerOptions}) do \texttt{CommonRequest}. \\
        \hline
    \end{tabularx}
    \caption{Mapeamento Inbound: Request da \textbf{Gemini} para o Domínio}
\end{table}

No sentido inverso, a resposta do domínio comum é convertida para o formato \texttt{Gemini\allowbreak{}Generate\allowbreak{}Content\allowbreak{}Response}.

\begin{table}[H]
    \centering
    \footnotesize
    \renewcommand{\arraystretch}{1.25}
    \begin{tabularx}{\textwidth}{|>{\ttfamily\raggedright\arraybackslash}p{4cm}|Y|}
        \hline
        \textbf{Domínio Comum} & \textbf{Mapeamento para Contrato Resposta (Gemini)} \\
        \hline
        id \& model & São atribuídos de forma direta aos campos \texttt{responseId} e \texttt{modelVersion} do \textbf{Gemini}. \\
        \hline
        choices & Dão origem aos respetivos objetos \texttt{GeminiCandidate}, preservando o índice da resposta. A conversão de papéis é feita no sentido inverso: \texttt{SYSTEM} e \texttt{ASSISTANT} passam a \texttt{"model"}, \texttt{TOOL} a \texttt{"function"}, e \texttt{USER} mantém-se como \texttt{"user"}. \\
        \hline
        choices[].message\newline.content & De forma simétrica, o conteúdo é traduzido consoante a sua natureza: um \texttt{TextPart} gera o texto base de um \texttt{GeminiResponsePart}; um \texttt{ToolCallPart} preenche o nó correspondente à chamada de função; um \texttt{JsonPart} é convertido numa \texttt{String} com o conteúdo JSON; e um \texttt{RefusalPart} é emitido como texto com o motivo da recusa. \\
        \hline
        choices[].finishReason & Mapeado para os motivos de paragem nativos do \textbf{Gemini}: \texttt{STOP}, \texttt{MAX\_TOKENS} ou \texttt{SAFETY}. No caso de \texttt{TOOL\_CALL}, a implementação emite \texttt{STOP}; valores \texttt{OTHER} ou nulos não são emitidos. \\
        \hline
        usage & Os contadores genéricos do domínio são convertidos e agrupados na estrutura \texttt{GeminiUsageMetadata}, alimentando os campos \texttt{promptTokenCount}, \texttt{candidatesTokenCount} e \texttt{totalTokenCount}. \\
        \hline
    \end{tabularx}
    \caption{Mapeamento Inbound: Domínio para \textbf{Gemini} Generate Content Response Unária}
\end{table}

    Durante o fluxo de \textit{streaming} via Server-Sent Events (\textbf{SSE}), a API \textbf{Gemini} diferencia-se de fornecedores como a \textbf{OpenAI} ou a \textbf{Anthropic} por não adotar objetos específicos para deltas parciais. Em vez disso, reutiliza a mesma estrutura da resposta unária (\texttt{Gemini\allowbreak{}Generate\allowbreak{}Content\allowbreak{}Response}) para encapsular cada evento da \textit{stream}. Como consequência, o tradutor mapeia os deltas granulares de domínio (\texttt{CommonResponseEvent}) reconstruindo a hierarquia do DTO nativo a cada evento emitido.

\begin{table}[H]
    \centering
    \footnotesize
    \renewcommand{\arraystretch}{1.25}
    \begin{tabularx}{\textwidth}{|>{\raggedright\arraybackslash}p{4cm}|Y|}
        \hline
        \textbf{Common Response Event} & \textbf{Mapeamento para Evento SSE (Gemini Chunk)} \\
        \hline
        \texttt{ResponseStarted} & Emite uma instância básica contendo apenas os metadados iniciais, como o \texttt{modelVersion} e o \texttt{responseId}. \\
        \hline
        \texttt{ChoiceStarted} & Inicializa as instâncias de \texttt{GeminiCandidate}, definindo apenas os atributos essenciais, como o índice (\texttt{index}) e o papel (\texttt{role}). \\
        \hline
        \texttt{TextDeltaEvent} & Encaminha o fragmento de texto (\textit{delta}) reconstruindo a estrutura aninhada da resposta: \texttt{candidates[] -> content -> parts[] -> GeminiResponsePart(text)}. \\
        \hline
        \texttt{ToolCallStartedEvent} & Gera um candidato que assinala o início da chamada de função através de um \texttt{GeminiResponsePart(functionCall)}, configurando o seu nome e inicializando os argumentos com um \texttt{JsonObject} vazio. \\
        \hline
        \texttt{ToolCallArguments\newline DeltaEvent} & Transmite os fragmentos incrementais dos argumentos utilizando uma chave provisória (\texttt{"partialJson"}), que transporta a \textit{string} bruta para facilitar a posterior reconstrução do \textit{payload}. \\
        \hline
        \texttt{ChoiceFinished} & Emite apenas o motivo de paragem (\texttt{finishReason}) correspondente ao término desta resposta parcial (\textit{choice}); no caso de \texttt{TOOL\_CALL}, o valor nativo emitido é \texttt{STOP}. \\
        \hline
        \texttt{UsageReported} & Emite os contadores de \textit{tokens} consumidos, encapsulados num objeto \texttt{GeminiUsageMetadata}. \\
        \hline
        \texttt{ResponseCompleted} & Sinaliza o fim da \textit{stream}, emitindo um evento com uma estrutura idêntica à do \texttt{ResponseStarted}. \\
        \hline
        \texttt{ResponseErrored} & Emite uma resposta com \texttt{GeminiPromptFeedback}. Quando o erro é marcado como recuperável, o \texttt{blockReason} recebe o prefixo \texttt{TRANSIENT\_ERROR}; nos restantes casos, recebe o prefixo \texttt{ERROR}. \\
        \hline
    \end{tabularx}
    \caption{Mapeamento Inbound: Eventos \textbf{SSE} \textbf{Gemini} (\texttt{fromDomainEvent})}
\end{table}

    \section{Tradutores de Saída (Outbound)}
    \subsection{\textbf{OpenAI} Outbound Translator}

    A implementação \texttt{OpenAiOutboundTranslator} adapta o \texttt{CommonRequest} ao formato de \textit{chat completions} da \textbf{OpenAI}. Os campos comuns são mapeados diretamente sempre que possível, e as opções próprias da \textbf{OpenAI} são lidas da \texttt{TypedMap} de \textit{provider}.

    \begin{table}[H]
        \centering
        \footnotesize
        \renewcommand{\arraystretch}{1.25}
        \begin{tabularx}{\textwidth}{|>{\ttfamily\raggedright\arraybackslash}p{4cm}|Y|}
            \hline
            \textbf{Domínio Comum} & \textbf{Mapeamento para Contrato Saída (OpenAI Request)} \\
            \hline
            model & Mapeado de forma direta. \\
            \hline
            messages[].content & Os blocos \texttt{TextPart} são unidos com quebras de linha. Os \texttt{ToolCallPart} são convertidos em \texttt{OpenAiToolCall(id, function(name, arguments))}, usando a representação textual dos argumentos. Quando a mensagem tem papel \texttt{TOOL}, o primeiro \texttt{ToolResultPart} preenche o \texttt{content} e o respetivo \texttt{toolCallId}. \\
            \hline
            config & \texttt{temperature}, \texttt{maxTokens} e \texttt{topP} são transferidos para o pedido nativo. O \texttt{maxTokens} é limitado com \texttt{coerceAtMost(4000)} e assume \texttt{1000} por defeito. As \texttt{stopSequences} são convertidas para \texttt{OpenAiStop}. \\
            \hline
            providerOptions & São lidas, quando existem, as opções específicas \texttt{frequencyPenalty}, \texttt{presencePenalty}, \texttt{n}, \texttt{stream}, \texttt{seed}, \texttt{user}, \texttt{logitBias}, \texttt{logProbs} e \texttt{topLogProbs}. \\
            \hline
            tools & Cada \texttt{CommonTool} é convertido para \texttt{OpenAiFunctionDefinition(name, description, parameters)} dentro de um \texttt{OpenAiTool}. \\
            \hline
            toolChoice & Traduz \texttt{Auto}, \texttt{None} e \texttt{Required} para os modos nativos \texttt{auto}, \texttt{none} e \texttt{required}. Nos casos específicos, força a função através de \texttt{OpenAiToolChoice.Function}. \\
            \hline
            jsonResponse & Quando \texttt{jsonResponse} está ativo, define \texttt{OpenAiResponseFormat(type = "json\_object")}. \\
            \hline
        \end{tabularx}
        \caption{Mapeamento Outbound: CommonRequest para \textbf{OpenAI} Request}
    \end{table}

    Na receção de uma resposta unária, o tradutor converte a resposta da \textbf{OpenAI} para o modelo comum e preserva metadados úteis do fornecedor.

    \begin{table}[H]
        \centering
        \footnotesize
        \renewcommand{\arraystretch}{1.25}
        \begin{tabularx}{\textwidth}{|>{\ttfamily\raggedright\arraybackslash}p{4cm}|Y|}
            \hline
            \textbf{Contrato Resposta (OpenAI)} & \textbf{Mapeamento para o Domínio} \\
            \hline
            id \& model & Mapeados diretamente para \texttt{CommonResponse}. Os campos \texttt{created} e \texttt{systemFingerprint}, quando presentes, são guardados em \texttt{providerOptions}. \\
            \hline
            choices[].message.role & Convertido para \texttt{CommonRole}. Se o papel vier omisso, é assumido \texttt{assistant}. \\
            \hline
            choices[].message.content & O conteúdo textual é convertido para \texttt{TextPart}. Se o texto aparentar conter um objeto ou lista JSON e o \textit{parse} for válido, é convertido para \texttt{JsonPart}. As \texttt{toolCalls} da mensagem ou do \texttt{delta} são convertidas para \texttt{ToolCallPart}. \\
            \hline
            choices[].finishReason & Converte \texttt{stop}, \texttt{length}, \texttt{tool\_calls} e \texttt{content\_filter} para os motivos comuns correspondentes; valores desconhecidos ficam como \texttt{OTHER}. \\
            \hline
            usage & Mapeia \texttt{promptTokens}, \texttt{completionTokens} e \texttt{totalTokens} para \texttt{inputTokens}, \texttt{outputTokens} e \texttt{totalTokens}. \\
            \hline
        \end{tabularx}
        \caption{Mapeamento Outbound: \textbf{OpenAI} Response Unária para Domínio}
    \end{table}

    Para respostas em \textit{streaming} via \textbf{SSE}, o tradutor mantém o \texttt{id} e o \texttt{model} entre eventos, porque nem todos os \textit{chunks} da \textbf{OpenAI} repetem esses metadados.

    \begin{table}[H]
        \centering
        \footnotesize
        \renewcommand{\arraystretch}{1.25}
        \begin{tabularx}{\textwidth}{|>{\raggedright\arraybackslash}p{4cm}|Y|}
            \hline
            \textbf{Evento SSE (OpenAI)} & \textbf{Mapeamento para Common Response Event} \\
            \hline
            Estado do fluxo (\texttt{runningFold}) & Mantém um \texttt{OpenAiEventContext} com o último \texttt{id} e \texttt{model} conhecidos, reutilizando-os nos eventos seguintes quando o \textit{chunk} não os inclui. \\
            \hline
            \texttt{choices.first().delta} & Um \texttt{delta.role} origina \texttt{ChoiceStarted}. Um \texttt{delta.content} origina \texttt{TextDeltaEvent}. \\
            \hline
            \texttt{delta.toolCalls} & Se os argumentos parciais têm conteúdo e o nome da função está ausente, o evento é traduzido para \texttt{ToolCallArgumentsDeltaEvent}. Caso contrário, inicia uma chamada de ferramenta com \texttt{ToolCallStartedEvent}. \\
            \hline
            \texttt{finishReason} & Quando a \textit{choice} inclui motivo de conclusão, é emitido \texttt{ChoiceFinished} com o \texttt{finishReason} convertido para o domínio comum. \\
            \hline
            \texttt{usage} & Quando chega num \textit{chunk} sem \textit{choices}, é convertido para \texttt{UsageReported}. \\
            \hline
            \texttt{Done} ou \texttt{Chunk} final & O fim explícito do fluxo, ou um objeto que já não corresponde a \texttt{chat.completion.chunk}, é traduzido para \texttt{ResponseCompleted}. \\
            \hline
            \texttt{Error} & Gera \texttt{ResponseErrored}. O erro só é marcado como \texttt{retryable} quando o tipo é \texttt{server\_error}, \texttt{rate\_limit\_error} ou \texttt{temporarily\_unavailable}. \\
            \hline
        \end{tabularx}
        \caption{Mapeamento Outbound: \textbf{OpenAI} Streaming para Domínio (\texttt{toDomainEvent})}
    \end{table}

    \subsection{\textbf{Anthropic} Outbound Translator}

    O \texttt{AnthropicOutboundTranslator} adapta o domínio comum à \textit{Messages API} da \textbf{Anthropic}. O \texttt{systemPrompt} é enviado no campo próprio \texttt{system}, enquanto mensagens, ferramentas e opções específicas são organizadas segundo o contrato nativo.

    \begin{table}[H]
        \centering
        \footnotesize
        \renewcommand{\arraystretch}{1.25}
        \begin{tabularx}{\textwidth}{|>{\ttfamily\raggedright\arraybackslash}p{4cm}|Y|}
            \hline
            \textbf{Domínio Comum} & \textbf{Mapeamento para Contrato Saída (Anthropic Request)} \\
            \hline
            system & O \texttt{systemPrompt} é convertido para \texttt{RawText} quando contém texto não vazio. \\
            \hline
            messages[].role & \texttt{USER} é convertido para \texttt{user}; \texttt{SYSTEM}, \texttt{ASSISTANT} e \texttt{TOOL} são enviados como \texttt{assistant}. \\
            \hline
            messages[].content & Uma mensagem com um único \texttt{TextPart} é enviada como \texttt{RawText}. Nos restantes casos, é criado um \texttt{ListContentBlock}: \texttt{TextPart} gera \texttt{Text}, \texttt{ToolCallPart} gera \texttt{ToolUse}, \texttt{ToolResultPart} gera \texttt{ToolResult} e \texttt{JsonPart} é enviado como texto JSON. \\
            \hline
            config & Mapeia \texttt{temperature}, \texttt{topP} e \texttt{stopSequences}. O \texttt{maxTokens} usa o valor configurado ou \texttt{1024} por defeito. \\
            \hline
            providerOptions & Inclui opções próprias da \textbf{Anthropic}, como \texttt{stream}, \texttt{topK}, \texttt{stopToken}, \texttt{thinking} e \texttt{metadata}. \\
            \hline
            tools & Cada \texttt{CommonTool} é convertido para \texttt{AnthropicToolDefinition(name, description, inputSchema)}. \\
            \hline
            toolChoice & Converte \texttt{Auto}, \texttt{None} e \texttt{Required} para \texttt{auto}, \texttt{none} e \texttt{any}. Uma escolha específica é enviada como \texttt{type = "tool"} com o nome da função. \\
            \hline
        \end{tabularx}
        \caption{Mapeamento Outbound: CommonRequest para \textbf{Anthropic} Request}
    \end{table}


Na receção de uma resposta unária da \textbf{Anthropic} (\texttt{Anthropic\allowbreak{}Message\allowbreak{}Response}), o tradutor cria uma única \texttt{CommonChoice} e converte cada bloco de conteúdo para a representação comum.

    \begin{table}[H]
        \centering
        \footnotesize
        \renewcommand{\arraystretch}{1.25}
        \begin{tabularx}{\textwidth}{|>{\ttfamily\raggedright\arraybackslash}p{4cm}|Y|}
            \hline
            \textbf{Contrato Resposta (Anthropic)} & \textbf{Mapeamento para o Domínio} \\
            \hline
            id \& model & Mapeados diretamente para \texttt{CommonResponse}; a resposta é representada numa \texttt{CommonChoice} com índice \texttt{0}. \\
            \hline
            content[].type == text & Convertido para \texttt{TextPart}. \\
            \hline
            content[].type == tool\_use & Convertido para \texttt{ToolCallPart}, preservando o \texttt{id}, o nome da função e as propriedades do \texttt{input}. \\
            \hline
            content[].type == thinking & Como este tradutor não tem um bloco próprio para \texttt{thinking} no domínio comum, o conteúdo é preservado como \texttt{RefusalPart(reason = thinking)}. \\
            \hline
            stopReason & Converte \texttt{end\_turn}, \texttt{max\_tokens}, \texttt{tool\_use} e \texttt{stop\_sequence} para os motivos comuns correspondentes. \\
            \hline
            usage & Mapeia \texttt{inputTokens} e \texttt{outputTokens}; o total é calculado a partir desses dois valores quando disponíveis. \\
            \hline
        \end{tabularx}
        \caption{Mapeamento Outbound: \textbf{Anthropic} Response Unária para Domínio}
    \end{table}

    Em \textit{streaming} \textbf{SSE}, a \textbf{Anthropic} já separa o fluxo em eventos tipados. O tradutor converte cada tipo de evento para o evento comum correspondente e mantém em contexto o \texttt{id} e o \texttt{model} recebidos no início da mensagem.

    \begin{table}[H]
        \centering
        \footnotesize
        \renewcommand{\arraystretch}{1.25}
        \begin{tabularx}{\textwidth}{|>{\raggedright\arraybackslash}p{4cm}|Y|}
            \hline
            \textbf{Evento SSE (Anthropic)} & \textbf{Mapeamento para Common Response Event} \\
            \hline
            Estado do fluxo (\texttt{runningFold}) & Mantém um \texttt{AnthropicEventContext} com o \texttt{id} e o \texttt{model} da mensagem, reutilizando-os nos eventos seguintes do mesmo fluxo. \\
            \hline
            \texttt{MessageStart} & Convertido para \texttt{ResponseStarted}, usando o \texttt{id} e o \texttt{model} da mensagem nativa. \\
            \hline
            \texttt{ContentBlockStart} & Um bloco \texttt{Text} origina \texttt{ChoiceStarted}. Um bloco \texttt{ToolUse} origina \texttt{ToolCallStartedEvent}. Um bloco \texttt{Thinking} também inicia uma \textit{choice} de assistente. \\
            \hline
            \texttt{ContentBlockDelta} & \texttt{TextDelta}, \texttt{SignatureDelta} e \texttt{ThinkingDelta} são convertidos para \texttt{TextDeltaEvent}. \texttt{InputJsonDelta} é convertido para \texttt{ToolCallArgumentsDeltaEvent}. \\
            \hline
            \texttt{ContentBlockStop} & Convertido para \texttt{ChoiceFinished} no índice do bloco recebido. \\
            \hline
            \texttt{MessageDelta} & Quando inclui \texttt{usage}, gera \texttt{UsageReported}; caso contrário, gera \texttt{ChoiceFinished} com o \texttt{stopReason} convertido. \\
            \hline
            \texttt{MessageStop} & Convertido para \texttt{ResponseCompleted}. \\
            \hline
            \texttt{Ping} & Mantém o fluxo ativo e é representado como \texttt{ResponseStarted}. \\
            \hline
            \texttt{Error} & Gera \texttt{ResponseErrored}. O campo \texttt{retryable} é marcado quando o tipo de erro contém \texttt{overloaded}. \\
            \hline
        \end{tabularx}
        \caption{Mapeamento Outbound: \textbf{Anthropic} Streaming para Domínio (\texttt{toDomainEvent})}
    \end{table}

    \subsection{\textbf{Gemini} Outbound Translator}

    O \texttt{GeminiOutboundTranslator} adapta o domínio comum ao formato \texttt{GenerateContent} da API \textbf{Gemini}. A implementação trata ainda diferenças próprias da plataforma, como a limpeza dos JSON \textit{schemas} de ferramentas e a forma como os eventos progressivos reutilizam a estrutura de resposta completa.

    \begin{table}[H]
        \centering
        \footnotesize
        \renewcommand{\arraystretch}{1.25}
        \begin{tabularx}{\textwidth}{|>{\ttfamily\raggedright\arraybackslash}p{4cm}|Y|}
            \hline
            \textbf{Domínio Comum} & \textbf{Mapeamento para Contrato Saída (Gemini Request)} \\
            \hline
            systemInstruction & O \texttt{systemPrompt} é convertido para \texttt{GeminiSystemInstruction}, contendo uma lista de \texttt{GeminiPart} com o texto do \textit{prompt} sistémico. \\
            \hline
            messages[].content & \texttt{TextPart} é enviado como texto. \texttt{ToolCallPart} gera \texttt{functionCall} e inclui \texttt{thoughtSignature = "skip\_thought\_signature\_validator"}. \texttt{ToolResultPart} gera \texttt{functionResponse}. \texttt{JsonPart} é enviado como texto JSON. \\
            \hline
            config & \texttt{stopSequences}, \texttt{temperature} e \texttt{topP} são colocados em \texttt{GeminiGenerationConfig}. As opções \texttt{topK}, \texttt{thinkingConfig}, \texttt{responseMimeType} e \texttt{responseJsonSchema} são lidas de \texttt{providerOptions}; quando \texttt{jsonResponse} está ativo, o \texttt{responseMimeType} passa a \texttt{application/json}. \\
            \hline
            tools \& cleanGeminiParameters & As ferramentas são agrupadas num \texttt{GeminiTool} com \texttt{functionDeclarations}. O \texttt{parametersSchema} é limpo recursivamente para remover chaves não suportadas (\texttt{\$schema}, \texttt{additionalProperties}, \texttt{propertyNames}, \texttt{title}, \texttt{default}, \texttt{\$id}, \texttt{exclusiveMinimum}); valores \texttt{const} são convertidos para \texttt{enum}. \\
            \hline
            toolChoice & Converte \texttt{Auto}, \texttt{None} e \texttt{Required} para os modos \texttt{AUTO}, \texttt{NONE} e \texttt{ANY}. Em escolhas específicas, usa \texttt{ANY} com \texttt{allowedFunctionNames}. \\
            \hline
        \end{tabularx}
        \caption{Mapeamento Outbound: CommonRequest para \textbf{Gemini} Request}
    \end{table}

    Na resposta unária \texttt{GeminiGenerateContentResponse}, o tradutor percorre os candidatos devolvidos e converte as partes reconhecidas para o domínio comum.

    \begin{table}[H]
        \centering
        \footnotesize
        \renewcommand{\arraystretch}{1.25}
        \begin{tabularx}{\textwidth}{|>{\ttfamily\raggedright\arraybackslash}p{4cm}|Y|}
            \hline
            \textbf{Contrato Resposta (Gemini)} & \textbf{Mapeamento para o Domínio} \\
            \hline
            id \& model & Usa \texttt{responseId} quando existe; se estiver ausente, gera um identificador local por UUID. O \texttt{modelVersion} é usado como modelo e, se vier ausente, fica vazio. \\
            \hline
            candidates[].parts & Partes de texto geram \texttt{TextPart}. Partes \texttt{functionCall} geram \texttt{ToolCallPart}, com \texttt{toolCallId} local \texttt{gemini-tool-call-0} e argumentos convertidos a partir de \texttt{args}. \\
            \hline
            candidates[].finishReason & Se a resposta contém chamadas de ferramenta, o motivo comum fica como \texttt{TOOL\_CALL}. Caso contrário, \texttt{STOP}, \texttt{MAX\_TOKENS} e \texttt{SAFETY} são convertidos para \texttt{STOP}, \texttt{LENGTH} e \texttt{CONTENT\_FILTER}. \\
            \hline
            promptFeedback & Quando existe, o motivo de bloqueio é preservado em \texttt{providerOptions} sob a chave \texttt{promptFeedback}. \\
            \hline
            usageMetadata & Mapeia \texttt{promptTokenCount}, \texttt{candidatesTokenCount} e \texttt{totalTokenCount} para a estrutura comum de utilização. \\
            \hline
        \end{tabularx}
        \caption{Mapeamento Outbound: \textbf{Gemini} Response Unária para Domínio}
    \end{table}

    Nos eventos progressivos da API \textbf{Gemini}, cada \textit{chunk} chega com a mesma estrutura de \texttt{GeminiGenerateContentResponse}. O tradutor analisa o conteúdo disponível em cada emissão e produz o evento comum correspondente.

    \begin{table}[H]
        \centering
        \footnotesize
        \renewcommand{\arraystretch}{1.25}
        \begin{tabularx}{\textwidth}{|>{\raggedright\arraybackslash}p{4cm}|Y|}
            \hline
            \textbf{Evento SSE (Gemini)} & \textbf{Mapeamento para Common Response Event} \\
            \hline
            Estado do fluxo (\texttt{runningFold}) & Mantém um \texttt{GeminiEventContext} com o último \texttt{id} e \texttt{model} conhecidos, para preencher eventos seguintes quando o \textit{chunk} não repete esses dados. \\
            \hline
            \texttt{promptFeedback} & Quando inclui \texttt{blockReason}, gera \texttt{ResponseErrored}. O erro é marcado como \texttt{retryable} quando o motivo contém \texttt{TRANSIENT}. \\
            \hline
            \texttt{candidates} vazio & Se existir \texttt{usageMetadata}, gera \texttt{UsageReported}; caso contrário, assinala o início do fluxo com \texttt{ResponseStarted}. \\
            \hline
            \texttt{finishReason} & Quando presente num candidato, gera \texttt{ChoiceFinished} com o motivo convertido para o domínio comum. \\
            \hline
            \texttt{functionCall} & Se o nome da função está presente, inicia a chamada com \texttt{ToolCallStartedEvent}. Como a \textbf{Gemini} não fornece identificador de ferramenta neste ponto, é criado um \texttt{toolCallId} local no formato \texttt{gemini-tool-call-\{Random\}}. Se o nome está ausente mas existem argumentos, emite \texttt{ToolCallArgumentsDeltaEvent}; a chave \texttt{partialJson} é usada quando existe, caso contrário é usada a representação textual dos argumentos. \\
            \hline
            \texttt{text} & Um fragmento textual em \texttt{parts} é convertido para \texttt{TextDeltaEvent}. \\
            \hline
            \texttt{content.role} & Quando o \textit{chunk} apenas indica o papel da resposta, é emitido \texttt{ChoiceStarted}. \\
            \hline
            \texttt{Done} & Convertido para \texttt{ResponseCompleted}. \\
            \hline
            \texttt{Error} & Gera \texttt{ResponseErrored}. O campo \texttt{retryable} é marcado para estados \texttt{UNAVAILABLE}, \texttt{RESOURCE\_EXHAUSTED}, \texttt{DEADLINE\_EXCEEDED}, \texttt{ABORTED} e \texttt{INTERNAL}. \\
            \hline
        \end{tabularx}
        \caption{Mapeamento Outbound: \textbf{Gemini} Streaming para Domínio (\texttt{toDomainEvent})}
    \end{table}

